\subsection{\texorpdfstring{Weakly closed invariant subspaces of \(L^{\infty}(G)\)}{Weakly closed invariant subspaces of L\^{}\{\textbackslash infinity\}(G)}}

In this issue, we identify \(\mathrm{L}^{\infty}(\mathrm{G})\) with the dual of \(\mathrm{L}^{1}(\mathrm{G})\) , and we endow it with the weak topology \(\sigma(\mathrm{L}^{\infty}(\mathrm{G}),\mathrm{L}^{1}(\mathrm{G}))\) . One denotes by \((f,g)\mapsto\langle f,g\rangle\) the bilinear map defining this duality for \(f\in\mathrm{L}^{1}(\mathrm{G})\) and \(g\in\mathrm{L}^{\infty}(\mathrm{G})\) .

The map W \(\mapsto\) W° is a bijection from the set of weakly closed vector subspaces of L∞(G) onto the set of closed vector subspaces of L¹(G) (EVT, II, p.~55, prop. 10).

On the other hand, if \(f \in \mathrm{L}^1(\mathrm{G})\) and \(x \in \mathrm{G}\) , the endomorphism \(g \mapsto f * g\) (resp. \(g \mapsto \varepsilon_x * g\) ) of the Banach space \(\mathrm{L}^1(\mathrm{G})\) has as its transpose the endomorphism \(h \mapsto \check{f} * h\) (resp. \(h \mapsto \varepsilon_{x^{-1}} * h\) ) of the Banach space \(\mathrm{L}^\infty(\mathrm{G})\) (INT, VIII, §4, no. 3, example 6). For a closed vector subspace of \(\mathrm{L}^1(\mathrm{G})\) to be an ideal of \(\mathrm{L}^1(\mathrm{G})\) , it is necessary and sufficient that it be invariant under translations of G. Therefore, for a weakly closed vector subspace of \(\mathrm{L}^\infty(\mathrm{G})\) to be stable under convolution with elements of \(\mathrm{L}^1(\mathrm{G})\) , it is necessary and sufficient that it be invariant under translations of G.

Let W be a weakly closed vector subspace of \(L^{\infty}(G)\) . Suppose W (and hence also \(W^{\circ}\) ) is invariant under translations of G. Let \(f \in L^{1}(G)\) . For every \(g \in L^{\infty}(G)\) , one has \((\check{f} * g)(x) = \langle \varepsilon_{x} * f, g \rangle = \langle f, \varepsilon_{x^{-1}} * g \rangle\) . Therefore, for f to belong to \(W^{\circ}\) , it is necessary and sufficient that \(\check{f} * g = 0\) for all \(g \in W\) .

If W is a weakly closed vector subspace of \(L^{\infty}(G)\) invariant under translations of G, we denote by A(W) the set of characters \(\chi \in \widehat{G}\) that belong to W. This is a closed subset of \(\widehat{G}\) . If F is a closed subset of \(\widehat{G}\) , we denote by Y(F) the weakly closed vector subspace of \(L^{\infty}(G)\) generated by the elements of F; since every translation of G transforms each character into a function proportional to that character, the space Y(F) is translation-invariant.

Let W be a weakly closed subspace of \(\mathrm{L}^{\infty}(\mathrm{G})\) invariant under translations of G. By the bipolar theorem (EVT, II, p.~48, th. 1), a character \(\chi\) belongs to W if and only if it belongs to \((\mathrm{W}^{\circ})^{\circ}\) ; this latter space is the set of functions \(g \in \mathrm{L}^{\infty}(\mathrm{G})\) such that \(\langle f, g \rangle = 0\) for every \(f \in W^{\circ}\) . We have \(\langle f, \chi \rangle = \overline{\mathcal{F}}(f)(\chi)\) , and hence

\[
\mathrm{A} (\mathrm{W}) = \mathrm{V} (\mathrm{W} ^ {\circ}).
\]

Similarly, a function \(f \in \mathrm{L}^1(\mathrm{G})\) belongs to \(\Upsilon(\mathrm{F})^\circ\) if and only \(\langle f, \chi \rangle = 0\) for all \(\chi \in \mathrm{F}\) , which is equivalent to \(\overline{\mathcal{F}}(f)(\chi) = 0\) for \(\chi \in \mathrm{F}\) , that is to say to \(f \in \Upsilon(\mathrm{F})\) . Therefore (loc. cit.) we have

\[
\mathrm{Y} (\mathrm{F}) = \Upsilon (\mathrm{F}) ^ {\circ}.
\]

The relations \(V(\Upsilon(F)) = F\) (I, p.~13 and I, p.~30) and \(\Upsilon(V(I)) \supset I\) , combined with the bipolar theorem (EVT, II, p.~48, th. 1), then imply

\[
\mathrm{A} (\mathrm{Y} (\mathrm{F})) = \mathrm{F}, \quad \mathrm{Y} (\mathrm{A} (\mathrm{W})) \subset \mathrm{W}.
\]

PROPOSITION 4. --- Let W be a nonzero weakly closed translation-invariant vector subspace of L∞(G). Then W contains at least one character of G.

We have seen that \(A(W) = V(W^{\circ})\) . Since \(W \neq 0\) , one has \(W^{\circ} \neq L^{1}(G)\) , and then \(V(W^{\circ})\) is nonempty by th. 1 of II, p.~252.

PROPOSITION 5. --- Let W be a weakly closed translation-invariant vector subspace of L∞(G).

\begin{enumerate}
\def\labelenumi{\alph{enumi})}
\item
  For any neighborhood U of A(W) in \(\widehat{G}\) , every function in W is a weak limit of linear combinations of characters belonging to U;
\item
  If the boundary of A(W) contains no nonempty perfect set, every function of W is a weak limit of linear combinations of characters belonging to W.
\end{enumerate}

To prove \(a\) ), it suffices by the bipolar theorem to show that if \(f\) is a function of \(\mathrm{L}^{1}(\mathrm{G})\) orthogonal to the elements of U, then \(f\) is orthogonal to W. Now, the Fourier cotransform \(\overline{\mathcal{F}}(f)\) then vanishes on the neighborhood U of \(\mathrm{A}(\mathrm{W}) = \mathrm{V}(\mathrm{W}^{\circ})\) , so that prop. 2 of II, p.~252 indeed shows that \(f \in \mathrm{W}^{\circ}\) . Assertion \(b\) ) is established in an analogous manner, using th. 2 of II, p.~257 instead of prop. 2 of II, p.~252.

